
\section{Exploratory Data Analysis}
\label{sec:eda}

This section presents an exploratory analysis of the dataset used in this thesis. The goal of this analysis is to understand the structure of the multimodal data before model development. Since the task combines chest X-ray images, clinical indication text, structured vital signs, and prior-study information, it is important to examine not only the label distribution but also the availability and quality of each input modality.

\subsection{Dataset Cohort and Subsets}
\label{subsec:eda_cohort}

The dataset used in this thesis is constructed from chest X-ray studies and associated clinical information. The image component is derived from MIMIC-CXR, while the clinical context is obtained from associated emergency department and hospital records. Each sample corresponds to a current chest X-ray study with available target labels and, when possible, linked clinical indication text, vital signs, and prior-study information.

The cohort construction process follows several filtering steps. First, chest X-ray studies are selected according to the target classification setting. Second, studies are matched with available clinical information, including indication text and structured vital signs. Third, prior-study information is identified for patients who have previous chest X-ray examinations before the current study. Cases with missing modalities are handled according to the preprocessing strategy described in Section~\ref{sec:methodology}.

Table~\ref{tab:cohort_summary} summarizes the final cohort used in this thesis, including the number of studies, patients, and available modalities across the training, validation, and test subsets.

\begin{table}[htbp]
    \centering
    \caption{Summary of the dataset cohort and train/validation/test splits.}
    \label{tab:cohort_summary}
    \begin{tabular}{lrrrrr}
        \hline
        Subset & Patients & Studies & Images & Text Available & Vitals Available \\
        \hline
        Train & 3,525 & 11,941 & 19,556 & 19,556 & 16,770 \\
        Validation & 542 & 1,804 & 3,047 & 3,047 & 2,571 \\
        Test & 1,000 & 3,333 & 5,642 & 5,642 & 4,868 \\
        Total & 5,067 & 17,078 & 28,245 & 28,245 & 24,209 \\
        \hline
    \end{tabular}
\end{table}

This summary is important because multimodal learning depends on the availability of each modality. If a large proportion of samples are missing clinical text, vital signs, or prior studies, the model must be designed to remain robust under incomplete information.

\subsection{Label Distribution and Co-occurrence}
\label{subsec:eda_labels}

The target task is multi-label chest X-ray classification, meaning that each study may contain zero, one, or multiple positive findings. This differs from standard single-label classification, where each sample belongs to only one class. In chest radiography, co-occurring findings are common because different thoracic abnormalities may appear together. For example, cardiomegaly and edema may co-occur in some cases, while pleural effusion may appear with atelectasis or consolidation.

The dataset also follows a long-tailed distribution. Some findings occur frequently, while others appear much less often. This creates a class imbalance problem because frequent labels contribute many more positive examples during training than rare labels. As a result, a model trained without imbalance-aware strategies may perform well on common findings but poorly on rare but clinically important abnormalities.

Figure~\ref{fig:label_distribution} shows the number of positive samples for each target label.

\begin{figure}[htbp]
    \centering
    \includegraphics[width=0.95\textwidth]{img/label_distribution.png}
    \caption{Distribution of positive labels across the target chest X-ray findings.}
    \label{fig:label_distribution}
\end{figure}

In addition to individual label frequencies, label co-occurrence is also important. Co-occurrence analysis helps identify which abnormalities frequently appear together and whether the dataset contains clinically meaningful label relationships. Figure~\ref{fig:label_cooccurrence} presents a co-occurrence matrix between target findings.

\begin{figure}[htbp]
    \centering
    \includegraphics[width=0.85\textwidth]{img/label_cooccurrence.png}
    \caption{Co-occurrence matrix of target pathology labels.}
    \label{fig:label_cooccurrence}
\end{figure}

These label statistics motivate the use of multi-label learning and imbalance-aware loss functions, which are described later in Section~\ref{sec:methodology}.

\subsection{Clinical Indication Text Analysis}
\label{subsec:eda_text}

Clinical indication text provides information about the reason for the chest X-ray examination. This may include symptoms, suspected diagnoses, or clinical concerns such as shortness of breath, fever, cough, trauma, chest pain, or follow-up evaluation. Unlike radiology reports, which are produced after image interpretation, clinical indication text is available before or during interpretation and is therefore suitable as input context for prediction.

The text modality has variable length and content. Some indications are short phrases, while others contain more detailed clinical descriptions. This variation affects the text encoder because very long sequences may need to be truncated, while very short sequences may contain limited information.

Figure~\ref{fig:text_length_distribution} shows the distribution of clinical indication text lengths.

\begin{figure}[htbp]
    \centering
    \includegraphics[width=0.85\textwidth]{img/text_length_distribution.png}
    \caption{Distribution of clinical indication text lengths (in words).}
    \label{fig:text_length_distribution}
\end{figure}

Table~\ref{tab:common_text_terms} lists frequent clinical terms found in the indication text. This helps verify that the text modality contains clinically relevant information rather than only administrative or generic phrases.

\begin{table}[htbp]
    \centering
    \caption{Most frequent clinical terms in the indication text.}
    \label{tab:common_text_terms}
    \begin{tabular}{lr}
        \hline
        Term or phrase & Frequency \\
        \hline
        old & 9,013 \\
        year & 8,967 \\
        man & 6,264 \\
        yearold & 4,771 \\
        woman & 4,628 \\
        \hline
    \end{tabular}
\end{table}

This analysis supports the use of a domain-specific text encoder in the multimodal model.

\subsection{Vital Signs and Missingness Analysis}
\label{subsec:eda_vitals}

Structured vital signs provide additional clinical context about the patient's physiological state near the time of imaging. In this thesis, the vital sign modality includes variables such as temperature, heart rate, oxygen saturation, respiratory rate, systolic blood pressure, and diastolic blood pressure. These values may help the model interpret radiographic findings in relation to the patient's current condition.

However, vital signs are not always complete. Missing values are common in real clinical datasets and must be handled carefully. Instead of simply discarding samples with missing vitals, this thesis uses a missingness-aware representation, where missing values are filled using a defined strategy and a corresponding missingness mask is provided to the model.

Table~\ref{tab:vitals_summary} summarizes the distribution and missingness of each vital sign.

\begin{table}[htbp]
    \centering
    \caption{Summary statistics and missingness rates for vital signs in the training split.}
    \label{tab:vitals_summary}
    \begin{tabular}{lrrrrrr}
        \hline
        Variable & Mean & Std & Median & Min & Max & Missing (\%) \\
        \hline
        Temperature & 98.15 & 3.37 & 98.20 & 1.04 & 105.40 & 17.18 \\
        Heart rate & 88.53 & 19.48 & 88.00 & 7.00 & 200.00 & 14.08 \\
        O2 saturation & 97.63 & 13.36 & 98.00 & 48.00 & 921.00 & 15.15 \\
        Respiratory rate & 18.64 & 3.60 & 18.00 & 1.00 & 65.00 & 15.46 \\
        Systolic BP & 130.94 & 24.65 & 129.00 & 8.00 & 581.00 & 14.27 \\
        Diastolic BP & 74.42 & 76.49 & 73.00 & 0.00 & 7,698.90 & 14.35 \\
        \hline
    \end{tabular}
\end{table}

Figure~\ref{fig:vitals_missingness} visualizes the missingness pattern across vital sign variables.

\begin{figure}[htbp]
    \centering
    \includegraphics[width=0.85\textwidth]{img/vitals_missingness.png}
    \caption{Missingness rate of each vital sign variable.}
    \label{fig:vitals_missingness}
\end{figure}

The presence of missing values motivates the use of both normalized vital values and missingness masks in the vital-sign encoder.

\subsection{Temporal Prior Gaps}
\label{subsec:eda_prior_gaps}

Prior-study information is available for patients who had previous chest X-ray examinations before the current study. In clinical practice, prior studies are useful because they allow comparison over time. A finding may be interpreted differently depending on whether it is new, stable, improving, or worsening. Therefore, the availability and timing of prior studies are important for prior-aware modelling.

Not all studies have an available prior examination. Among studies with priors, the time gap between the current and prior study may vary from a short interval to a much longer period. This creates two modelling challenges. First, the model must handle cases where no prior is available. Second, the model must understand that a recent prior may be more informative than a very old prior.

Table~\ref{tab:prior_availability} summarizes the availability of prior studies in the dataset.

\begin{table}[htbp]
    \centering
    \caption{Availability of prior studies at study level.}
    \label{tab:prior_availability}
    \begin{tabular}{lrr}
        \hline
        Prior-study status & Count & Percentage \\
        \hline
        With prior study & 14,045 & 82.24\% \\
        Without prior study & 3,033 & 17.76\% \\
        Total & 17,078 & 100.00\% \\
        \hline
    \end{tabular}
\end{table}

Figure~\ref{fig:prior_gap_distribution} shows the distribution of time gaps between current and prior studies.

\begin{figure}[htbp]
    \centering
    \includegraphics[width=0.85\textwidth]{img/prior_gap_distribution.png}
    \caption{Distribution of time gaps between current and prior chest X-ray studies (in days, log-scaled).}
    \label{fig:prior_gap_distribution}
\end{figure}

Because the time-gap distribution may be skewed, this thesis discretizes the time difference into temporal buckets before adding it to the model as a learned embedding. This allows the model to represent approximate temporal distance without assuming that time differences have a simple linear effect.
